\section{Phương pháp và thiết kế thực nghiệm}
\label{sec:methodology}

\subsection{Các mã cấu hình thí nghiệm}

Bảy mã E0--E6 được định nghĩa trong giao thức; sáu cấu hình hoạt động được đưa vào chiến dịch so sánh cuối cùng. Ký hiệu E chỉ là mã cấu hình thí nghiệm để rút gọn bảng và truy nguyên artifact, không biểu thị thứ hạng hay phiên bản mô hình.

\begin{table}[H]
\centering
\caption{Các mã cấu hình, mục tiêu và trạng thái phân tích}
\label{tab:experiment_matrix}
\begin{tabularx}{\textwidth}{lYYY}
\toprule
\textbf{Mã} & \textbf{Bộ trích đặc trưng/tiền xử lý} & \textbf{Mô hình chuỗi} & \textbf{Mục tiêu hoặc trạng thái} \\
\midrule
E0 & CNN của mốc tham chiếu & BiLSTM & Đối chứng nội bộ \\
E1 & Giữ encoder/đặc trưng E0 & TCN & Thay cấu hình mô hình chuỗi và lịch huấn luyện \\
E2 & ResNet-1D dùng epoch hiện tại & TCN & Thay gói đặc trưng/ngữ cảnh E1 \\
E3 & ResNet-1D và xử lý biên độ chính & TCN & Cấu hình đề xuất \\
E4 & ResNet-1D và lọc dải & TCN & Đối chiếu tiền xử lý \\
E5 & ResNet-1D, lọc dải và cắt $\pm800~\mu\mathrm{V}$ & TCN & Kiểm toán đầu vào đồng nhất với E4 \\
E6 & ResNet-1D và z-score theo bản ghi & TCN & Đối chiếu đổi thang \\
\bottomrule
\end{tabularx}
\end{table}

E0 là đối chứng nội bộ đã hiệu chỉnh theo giao thức v2, không phải bản sao định lượng của bài báo gốc. Vì vậy, các đối chiếu trong nghiên cứu được hiểu là so sánh giữa các cấu hình trong cùng thí nghiệm.

E1 dùng lại đúng encoder và cache đặc trưng E0, thay cấu hình BiLSTM bằng TCN cùng lịch huấn luyện tương ứng. Learning rate, batch size và quy tắc dừng sớm của hai mô hình chuỗi khác nhau; E1--E0 vì vậy không cô lập tác động kiến trúc. E1 cung cấp vector 75 chiều từ các nhánh C/P/N, còn E2 dùng embedding ResNet-1D 128 chiều chỉ từ epoch hiện tại. E2--E1 đồng thời thay encoder, ngữ cảnh P/N tường minh, số chiều đặc trưng và cách huấn luyện/chọn encoder; đây là đối chiếu gói đặc trưng/ngữ cảnh, không phải phép cô lập riêng encoder.

Encoder E0/E1 gồm 15 CNN, mỗi CNN có đầu ra softmax năm lớp. Các nhóm C/P/N dùng epoch hiện tại, liền trước hoặc liền sau để dự đoán nhãn của epoch trung tâm; dịch ngữ cảnh chỉ diễn ra trong từng bản ghi, lặp epoch đầu/cuối tại biên. Ghép các đầu ra tạo vector 75 chiều, không mặc định là xác suất đã hiệu chuẩn. ResNet-1D dùng ba khối dư 64/128/128 kênh và pooling toàn cục để tạo embedding 128 chiều. TCN gồm sáu khối, hai tích chập kernel 3 mỗi khối, dilation 1/2/4/8/16/32 và dropout 0,2. Padding đối xứng cho trường tiếp nhận lý thuyết 253 epoch. Mô hình chuỗi xử lý toàn bản ghi với padding và mặt nạ hợp lệ; cửa sổ 100 epoch chỉ thuộc phép benchmark, không phải cửa sổ huấn luyện.

E5 ban đầu được dùng để tách tác động của phép cắt biên $\pm800~\mu\mathrm{V}$ khỏi lọc dải E4. Kiểm toán toàn bộ 153 bản ghi cho thấy E4 và E5 giống nhau từng bit và tỷ lệ mẫu bị cắt bằng 0; do đó E5 không tạo điều kiện dữ liệu độc lập. Fold 00 được giữ làm bằng chứng kiểm toán, còn E5 bị loại trước các bảng hiệu năng và kiểm định thống kê để tránh một đối chiếu giả tạo giữa hai đầu vào đồng nhất.

\subsection{Bản đồ câu hỏi, bằng chứng và quyết định}

Không phải mọi Gate đều được xem là một đóng góp khoa học độc lập. Chuỗi Gate được dùng để bảo vệ đường đi từ dữ liệu đến kết luận, còn quyết định phát triển chỉ được rút ra khi có đại lượng và đối chứng phù hợp.

\begin{table}[H]
\centering
\caption{Câu hỏi thực tế và bằng chứng dùng để trả lời}
\label{tab:question_evidence_map}
\begin{tabularx}{\textwidth}{YYY}
\toprule
\textbf{Quyết định} & \textbf{Bằng chứng chính} & \textbf{Đầu ra được phép kết luận} \\
\midrule
Có nên thay cấu hình/gói đặc trưng? & Gate 5, Gate 6; E1--E0 và E2--E1 & Hiệu ứng cấu hình và lịch huấn luyện E1; hiệu ứng gói E2; đánh đổi tốc độ--tài nguyên \\
Pipeline nào đáng mang sang miền mới? & E3--E6 trong miền; E3--E0/E6 trên SHHS1; E3--E2 và E4 mở rộng & Thứ hạng toàn pipeline; không quy nhân quả cho một thao tác \\
Lỗi chuyển miền tập trung ở đâu? & F1 theo lớp, tỷ lệ dự đoán/nhãn thật, ma trận nhầm lẫn, vùng thay đổi nhãn & Xác định kênh lỗi đáng kiểm chứng tiếp; không suy ra cơ chế sinh lý \\
Các phân tích phụ có thay đổi quyết định không? & Silhouette và Gate 8 & Giới hạn của proxy hình học và ngữ cảnh mã hóa trùng lặp \\
Kết quả có truy nguyên được không? & Gate 1--4, Gate 7 và manifest & Tính hợp lệ của artifact; không phải bằng chứng về độ chính xác \\
\bottomrule
\end{tabularx}
\end{table}

\subsection{Huấn luyện và khóa đánh giá}

Các cấu hình dùng cùng các fold theo đối tượng. Encoder CNN15 được chọn theo weighted validation loss; ResNet-1D và mô hình chuỗi được chọn theo validation Macro-F1. Encoder đã chọn được đóng băng trước khi tạo cache và huấn luyện mô hình chuỗi; E1 tái sử dụng encoder/cache của E0. Tập test được giữ kín trong giai đoạn lựa chọn mô hình. Mỗi đối tượng chỉ có một vai trò test ngoài fold trong mỗi cấu hình và seed.

\begin{table}[H]
\centering
\caption{Thiết lập huấn luyện theo thành phần; patience tính theo số lần validation không cải thiện}
\label{tab:component_training}
\small
\begin{tabularx}{\textwidth}{lYYYY}
\toprule
 & \textbf{CNN15} & \textbf{ResNet-1D} & \textbf{BiLSTM} & \textbf{TCN} \\
\midrule
Optimizer & Adam & AdamW & Adam & Adam \\
Learning rate & 0,001 & 0,001 & 0,01 & 0,0005 \\
Weight decay & 0 & 0,0001 & 0 & 0 \\
Batch size & 64 epoch & 64 epoch & 4 bản ghi & 8 bản ghi \\
Epoch tối đa & 1.000 & 40 & 1.000 & 300 \\
Patience & 10 & 8 & 10 & 30 \\
Chọn checkpoint & Weighted loss & Macro-F1 & Macro-F1 & Macro-F1 \\
\bottomrule
\end{tabularx}
\end{table}

Mỗi CNN chuyên biệt lớp $c$ vẫn phân loại năm lớp với trọng số cross-entropy $w_c=0,5/(n_c/N)$ và $w_{k\ne c}=0,5/((N-n_c)/N)$. ResNet dùng $w_c=N/(5n_c)$; các số đếm $n_c$ và $N$ chỉ lấy từ tập huấn luyện. Mô hình chuỗi dùng masked cross-entropy với trọng số lớp bằng nhau và gradient clipping norm 1,0. Validation được thực hiện cuối mỗi training epoch. Nhãn không hợp lệ và padding bị loại khỏi loss và chỉ số. Khác biệt ở Bảng~\ref{tab:component_training} là một phần của cấu hình được so sánh, không được gộp vào một tác động kiến trúc riêng.

Quy trình này tách lựa chọn mô hình khỏi đánh giá cuối và bảo đảm các cấu hình sử dụng cùng khóa thời gian, cùng nhãn và cùng đơn vị suy luận. Mô hình đã chọn, dự đoán và chỉ số được kiểm tra tính nhất quán trước khi tổng hợp.

\subsection{Quy trình kiểm soát và khóa đánh giá (mã nội bộ Gate 1--8)}

Các cổng được tổ chức theo chuỗi bằng chứng:

\begin{enumerate}
  \item Gate 1 kiểm tra nguồn dữ liệu, nhãn, tiền xử lý và phép chia theo đối tượng.
  \item Gate 2 xác nhận khả năng chạy và phục hồi của quy trình.
  \item Gate 3 hoàn tất lựa chọn mô hình trên validation mà không sử dụng test.
  \item Gate 4 mở test theo quy tắc đã khóa và kiểm tra tính đầy đủ của dự đoán.
  \item Gate 5 thực hiện các so sánh hiệu năng bắt cặp.
  \item Gate 6 đánh giá độ phức tạp, tốc độ và hình học biểu diễn; chỉ hai đại lượng đầu trực tiếp trả lời quyết định vận hành.
  \item Gate 7 tập hợp bảng, hình và ma trận bằng chứng.
  \item Gate 8 khảo sát bổ sung giá trị dự báo biên của các nhóm ngữ cảnh mã hóa liền trước/liền sau tại vùng chuyển pha.
\end{enumerate}

\subsection{Phân tích thống kê}

Macro-F1 gộp trên Sleep-EDF được tính từ tổng các ma trận nhầm lẫn test ngoài fold, không phải trung bình có trọng số số epoch của Macro-F1 từng người. Khoảng tin cậy 95\% của chênh lệch gộp dùng bootstrap bắt cặp theo cụm đối tượng: các đêm của một người và các cấu hình được lấy lại mẫu cùng nhau. Wilcoxon signed-rank hai phía dùng chênh lệch Macro-F1 của từng đối tượng; đây không phải kiểm định chênh lệch trung bình. Macro-F1 luôn lấy trung bình trên năm lớp cố định, với F1 bằng 0 khi mẫu số của lớp bằng 0. Hai phép phân tích mô tả những khía cạnh khác nhau; bất đồng giữa chúng không tự xác định nhóm người nào tạo ra chênh lệch gộp. CI là khoảng marginal 95\%, còn Holm hiệu chỉnh các $p$-value trong bốn đối chiếu định trước của từng seed: E1--E0 (cấu hình mô hình chuỗi và lịch huấn luyện), E2--E1 (gói đặc trưng/ngữ cảnh), E3--E2 (chế độ xử lý chính) và E3--E6 (đổi thang so với z-score)~\cite{efron1994bootstrap,wilcoxon1945,holm1979}.

Đối chiếu toàn bộ cấu hình với mốc ban đầu được thực hiện sau khi đã quan sát kết quả nên được ghi rõ là hậu nghiệm. Các phân tích bổ sung không được nhập ngược vào nhóm giả thuyết chính.

\subsection{Độ nhạy theo seed}

Toàn bộ chiến dịch được lặp lại với một seed huấn luyện thứ hai trên cùng phép chia và cùng cấu hình. Hai lần chạy được báo cáo riêng, không gộp $p$-value và không xem hai seed cố định là mẫu đại diện cho mọi khởi tạo. Mục tiêu của lần lặp là đánh giá độ bền của hướng hiệu ứng, không phải xây dựng một ước lượng đầy đủ về phân phối hiệu năng.

\subsection{Benchmark và không gian đặc trưng}

Benchmark được thực hiện trên cùng phần cứng cho tất cả cấu hình. Phép đo phản ánh suy luận forward và không bao gồm đọc dữ liệu, tiền xử lý, lưu trữ trung gian hoặc huấn luyện. Thời gian wall-clock của chiến dịch seed 123 được báo cáo riêng, bao gồm huấn luyện, validation và hoàn thiện artifact trên 10 fold; các lượt được chạy tuần tự và test chưa được mở. Phân tích không gian đặc trưng so sánh các biểu diễn trên cùng epoch và cùng fold; Silhouette được dùng như chỉ số hỗ trợ~\cite{rousseeuw1987silhouettes}, còn t-SNE chỉ dùng để minh họa.

\subsection{Ablation nhóm ngữ cảnh (Gate 8 nội bộ)}

Gate 8 giữ họ TCN và lịch huấn luyện, thay đổi các nhóm đặc trưng C/P/N được cung cấp. Nhóm bị loại được thay bằng trung bình theo từng chiều từ các epoch có nhãn hợp lệ của tập huấn luyện từng fold. Cùng phép thay áp dụng cho train, validation và test, giữ nguyên 75 chiều; TCN được huấn luyện lại cho từng điều kiện CP, CN và C. Điều kiện đầy đủ CPN tái sử dụng E1. Đây không phải phép xóa đặc trưng chỉ lúc suy luận; validation và test không tham gia ước lượng giá trị thay thế.

Tiêu chí chính là Macro-F1 tại vùng chuyển pha. Ba so sánh ablation dùng cùng bootstrap cụm, Wilcoxon theo đối tượng và hiệu chỉnh Holm. Gate 8 đo hiệu ứng dự báo có điều kiện trong một quy trình cụ thể; không đo tỷ lệ thông tin, không xác định quan hệ nhân quả và không thiết lập tương đương.

\subsection{Đánh giá chuyển miền không cập nhật trọng số trên SHHS1}

Chiến dịch SHHS được tiến hành sau khi các mô hình Sleep-EDF đã được khóa, không cập nhật trọng số bằng SHHS. Với mỗi cấu hình, xác suất của mười mô hình outer fold được lấy trung bình cho từng epoch rồi chọn lớp có xác suất lớn nhất. Trên Sleep-EDF, mỗi người chỉ được dự đoán bởi mô hình ngoài fold chưa huấn luyện trên người đó. Vì cách tổ hợp mô hình và thành phần cohort khác nhau, chênh lệch điểm giữa Sleep-EDF và SHHS1 là chênh lệch mô tả giữa hai điều kiện đánh giá, không phải ước lượng riêng tác động dịch chuyển miền. Chỉ số chính SHHS là trung bình Macro-F1 từng người; bootstrap theo đối tượng và Wilcoxon bắt cặp được báo cáo với family hai đối chiếu chính E3--E0/E6, tách khỏi các family mở rộng.

Riêng E6 tính trung bình và độ lệch chuẩn trên toàn bộ tín hiệu của từng bản ghi SHHS đích sau lọc và cắt biên, trước cắt cửa sổ ngủ, không dùng nhãn. E6 do đó có phụ thuộc transductive ở cấp bản ghi; E0/E1/E2/E3 không ước lượng thống kê chuẩn hóa từ bản ghi đích. Thuật ngữ ``zero-shot'' trong tên chiến dịch chỉ việc không huấn luyện/cập nhật trọng số bằng SHHS. Nhãn tham chiếu vẫn được dùng để dựng cửa sổ đánh giá, xác định vùng nhãn và tính chỉ số; tất cả mô hình trong nghiên cứu đều được đánh giá ngoại tuyến.

Kết quả SHHS không được nhập ngược vào Gate 1--8. Đây là đánh giá ngoài miền trên mẫu đã khóa, nhằm kiểm tra độ bền của hướng hiệu ứng chứ không phải kiểm định lâm sàng hoặc đánh giá trên toàn bộ cohort SHHS.

\subsection{Định vị lỗi chuyển miền và phân tích phản thực}

Để xác định cấu trúc lỗi ngoài chỉ số tổng hợp, báo cáo phân tích precision, recall, F1, tỷ lệ số dự đoán trên số nhãn thật và ma trận nhầm lẫn theo lớp. Tỷ lệ dự đoán/nhãn thật mô tả tần suất phát lớp, không đo calibration xác suất. E6 kiểm tra liệu pipeline z-score đã huấn luyện tương ứng có khắc phục thiếu N3 hay không; đối chiếu này không phải can thiệp riêng lên biên độ của cùng một mô hình đóng băng.

Với $j$ là epoch đầu mang nhãn mới sau một cặp nhãn hợp lệ liên tiếp, vùng lân cận là hợp của các tập $\{j-1,j,j+1\}$. Khoảng trống nhãn không tạo một thay đổi liên tục. Vùng ổn định gồm các epoch hợp lệ cách ít nhất ba epoch so với cả hai vị trí $j-1$ và $j$ của mọi thay đổi nhãn trong bản ghi; bản ghi không có thay đổi nhãn cũng được xử lý theo quy tắc này. Hai vùng không bù nhau: 21.539 epoch Sleep-EDF và 21.073 epoch SHHS1 nằm ngoài cả hai vùng, không được gộp vào vùng ổn định. Kiểm tra persistent chỉ giữ thay đổi có ít nhất ba epoch liên tiếp cùng nhãn ở mỗi phía. Các vùng đều được xác định từ nhãn tham chiếu, nên là chẩn đoán ngoại tuyến chứ không phải bộ phát hiện chuyển pha sinh lý.

Phân tích phản thực (oracle) chuyển các số đếm ở đúng ô nhầm lẫn đã chọn sang ô dự đoán đúng trong ma trận E3 rồi tính lại \mf; nhãn tham chiếu và hỗ trợ từng lớp không đổi. Phép tính định lượng quy mô của các kênh lỗi trong cùng ma trận, không dự báo hiệu năng một thuật toán có thể đạt hoặc hiệu quả một chiến lược thu nhãn. Tỷ phần oracle dùng chênh lệch điểm quan sát giữa EDF ngoài fold và SHHS ensemble làm mẫu số. Các phần riêng không cộng tuyến tính do Macro-F1 phi tuyến; kết quả đồng thời có thể vượt chênh lệch giữa hai điểm gốc.
